To comprehensively assess the general visual question answering (VQA) capabilities of the Qwen3-VL series, we conduct extensive evaluations on a diverse set of benchmarks, including MMBench-V1.1~\citep{MMBench}, RealWorldQA~\citep{realworldqa2024}, MMStar~\citep{chen2024we}, and SimpleVQA~\citep{cheng2025simplevqa}. As detailed in~\cref{tab:results_flagship},~\cref{tab:results_medium} and~\cref{tab:results_small}, the Qwen3-VL family demonstrates robust and highly competitive performance across a wide spectrum of model sizes, from 2B to 235B parameters.

In the comparison of thinking mode, \ThinkBig achieves the highest score of 78.7 on MMStar. Gemini-2.5-Pro's~\citep{comanici2025gemini} Thinking mode delivers the best overall performance, but \ThinkBig is not far behind. In the non-reasoning mode comparison, \InstructBig obtains the highest scores on MMBench and RealWorldQA, with 89.3/88.9 and 79.2, respectively. 

In the experiments with medium-sized models, Qwen3-VL-32B-Thinking achieves the highest scores on MMBench and RealWorldQA, with 89.5/89.5 and 79.4, respectively. Notably, Qwen3-VL-32B-Instruct even outperforms the Thinking variant on RealWorldQA, scoring 79.0.

The scalability of the Qwen3-VL series is evident in the strong performance of our smaller models. Specifically, the largest model, Qwen3-VL-8B, achieves the highest performance across all five benchmarks. For example, on MMBench-EN, the score in "thinking" mode increases from 79.9 for the 2B model to 85.3 for the 8B model. A similar upward trend is observed on other benchmarks, such as MMStar, where the score rises from 68.1 (2B, thinking) to 75.3 (8B, thinking).


\begin{comment}
\begin{table}[t]
\centering
\caption{\textbf{Performance of Qwen3-VL-235BA22B and previous models on General VQA benchmarks.}}
\label{tab:vqa_results}
% \small
\setlength{\tabcolsep}{3.8pt}
\begin{tabular}{lll|llllll} \toprule
     \multirow[t]{3}{*}{\textbf{Benchmark}} & \textbf{Qwen3-VL} & \textbf{Qwen3-VL} & \textbf{Gemini} & \textbf{Gemini} & \textbf{OpenAI}  & \textbf{OpenAI}  & \textbf{Claude} & \textbf{Claude}  \\
    &  \textbf{235BA22B} & \textbf{235BA22B} & \textbf{2.5 Pro}& \textbf{2.5 Pro} & \textbf{GPT5}  & \textbf{GPT5} & \textbf{opus 4.1} & \textbf{opus 4.1}  \\
    &  {\scriptsize thinking} & {\scriptsize instruct} & {\scriptsize thinking} & {\scriptsize budget-128} & {\scriptsize thinking} & {\scriptsize no-thinking}& {\scriptsize thinking} & {\scriptsize no-thinking}  \\ 
\midrule
MMBench-V1.1-EN\textsubscript{\scriptsize test} & 88.8 & \underline{89.3} & $\mathbf{90.1^*}$ & 86.6 & 85.9 & 81.4 & 84.9 & 84.1 \\ 
MMBench-V1.1-CN\textsubscript{\scriptsize test} & 88.6 & \underline{88.9} & $\mathbf{89.7^*}$ & 86.4 & 83.5 & 79.9 & 84.9 & 74.3 \\ 
RealWorldQA & 81.3 & \underline{79.2} & $78.0^*$ & 76.0 & \textbf{82.8} & 73.3 & 69.9 & 68.5 \\ 
MMStar & \textbf{78.7} & 78.4 & $77.5^*$ & \underline{78.5} & 76.4 & 65.2 & 72.1 & 71.0 \\ 
SimpleVQA & 61.3 & 63.0 & \textbf{65.4} & \underline{66.9} & 61.8 & 56.7 & 56.7 & 55.7 \\ 
\midrule
Average & 79.7 & \underline{79.8} & \textbf{80.1} & 78.9 & 78.1 & 71.3 & 73.7 & 70.7 \\ 
\bottomrule
\end{tabular}
\end{table}


\begin{table}[!h]
\centering
\caption{\textbf{Performance of Qwen3-VL-32B, Qwen3-VL-30BA3B on General VQA benchmarks.}}
\label{tab:vqa_results_32_30}
% \small
\begin{adjustbox}{max width=1\textwidth}
% Note: The column specifier was changed from lll|ll|llll to l ll|ll|llll to match the new order.
\begin{tabular}{l ll|ll|llll} \toprule
     \multirow[t]{3}{*}{\textbf{Benchmark}} & \textbf{Qwen3-VL} & \textbf{Qwen3-VL} & \textbf{Qwen3-VL} & \textbf{Qwen3-VL} & \textbf{Gemini2.5}  & \textbf{Gemini2.5}  & \textbf{Qwen2.5-VL}   \\
    & \textbf{32B}& \textbf{32B} &  \textbf{30BA3B} & \textbf{30BA3B} & \textbf{flash}  & \textbf{flash} & \textbf{72B}  \\
    & {\scriptsize thinking} & {\scriptsize instruct} &  {\scriptsize thinking} & {\scriptsize instruct} & {\scriptsize thinking} & {\scriptsize nothinking}& {\scriptsize nothinking}  \\ 
\midrule
MMBench-V1.1-EN\textsubscript{\scriptsize test} & \textbf{89.5} & 87.6 & 87.0 & 86.1 & 87.1 & 86.6 & $\underline{88.4^*}$ \\ 
MMBench-V1.1-CN\textsubscript{\scriptsize test} & \textbf{89.4} & \underline{87.7} & 85.9 & 85.3 & 87.3 & 86.0 & 87.3 \\ 
RealWorldQA & \textbf{78.4} & \underline{79.0} & 77.4 & 73.7 & 76.0 & 75.7 & $75.7^*$ \\ 
MMStar & \textbf{79.4} & \underline{77.7} & 75.5 & 72.1 & 76.5 & 75.8 & $70.8^*$ \\ 
SimpleVQA & 55.4 & 56.9 & 54.3 & 52.7 & \textbf{63.2} & \underline{59.2} & 58.2 \\ 
\midrule
Average & \textbf{78.4} & \underline{77.8} & 76.0 & 74.0 & 78.0 & 76.7 & 76.1 \\ 
\bottomrule
\end{tabular}
\end{adjustbox}
\end{table}

\begin{table}[!h]
\centering
\caption{\textbf{Performance of Qwen3-VL-8B, Qwen3-VL-4B, Qwen3-VL-2B and previous models on General VQA benchmarks.}}
\label{tab:vqa_results_8_4_2_nano}
% \small
\begin{adjustbox}{max width=1\textwidth}
% Note: The column specifier was lll|ll|ll|ll, corrected to l ll|ll|ll|ll to match column count
\begin{tabular}{l ll|ll|ll|ll} \toprule
     \multirow[t]{3}{*}{\textbf{Benchmark}} & \textbf{Qwen3-VL} & \textbf{Qwen3-VL} & \textbf{Qwen3-VL} & \textbf{Qwen3-VL} & \textbf{Qwen3-VL}  & \textbf{Qwen3-VL}  & \textbf{GPT-5} & \textbf{GPT-5}  \\
    & \textbf{8B}  & \textbf{8B} & \textbf{4B}& \textbf{4B} & \textbf{2B} & \textbf{2B} & \textbf{Nano} & \textbf{Nano}  \\
    & {\scriptsize thinking} & {\scriptsize instruct}& {\scriptsize thinking} & {\scriptsize instruct} &  {\scriptsize thinking} & {\scriptsize instruct} & {\scriptsize high} & {\scriptsize minimal}  \\ 
\midrule
MMBench-V1.1-EN\textsubscript{\scriptsize test} & \textbf{85.3} & 84.5 & \underline{84.6} & 83.9 & 79.9 & 78.4 & 78.4 & 50.8 \\ 
MMBench-V1.1-CN\textsubscript{\scriptsize test} & \textbf{85.5} & \underline{84.7} & 83.8 & 83.5 & 78.4 & 75.9 & 77.6 & 48.5 \\ 
RealWorldQA & \textbf{73.5} & \underline{71.5} & 73.2 & 70.9 & 69.5 & 63.9 & 71.8 & 60.9 \\ 
MMStar & \textbf{75.3} & \underline{70.9} & 73.2 & 69.8 & 68.1 & 58.3 & 68.6 & 41.3 \\ 
SimpleVQA & \textbf{49.6} & \underline{50.2} & 48.8 & 48.0 & 43.6 & 40.7 & 46.0 & 39.0 \\
\midrule
Average & \textbf{73.8} & \underline{72.4} & 72.7 & 71.2 & 67.9 & 63.4 & 68.5 & 48.1 \\
\bottomrule
\end{tabular}
\end{adjustbox}
\end{table}
\end{comment}
% \subsection{Subjective Experience and Instruction Following}
% The ability to follow complex user instructions and reduce potential image-level hallucinations is indispensable for current large vision language models (VLMs). Here we assess our models on three such representative benchmarks: MM-MT-Bench~\citep{mm_mt_bench}, HallusionBench~\citep{hallusion_bench} and MIA-Bench~\citep{mia-bench}. MM-MT-Bench is a multi-turn LLM-as-a-judge evaluation benchmark for testing multimodal instruction-tuned models. HallusionBench aims at diagnosing image-context reasoning and poses great challenges for current VLMs. MIA-Bench is a more comprehensive benchmark to evaluate models' reactions to user's complex instructions (e.g., creative writing with character limit and compositional instructions).